\chapter{第四章补充材料}

\section{术语和符号}\label{sec-list-of-terms-and-notations}
本节包含术语和符号列表. 
\begin{itemize}
    \item 粗体字母 (例如, $\x$, $\v$ 和 $\w$) 表示向量. $\x$ 的第 $i$ 个分量记为 $\x^{(i)}$. 
    \item $[N] := \{1, 2, \cdots, N\}$ 表示从 $1$ 到 $N$ 的指标集. 
    \item $[\x_{N}] = \{\x : \exists i \in [N] \ s.t.\ \x = \x_i\}$  (如果 $i \neq j \implies \x_i \neq \x_j$, 则 $ = \{\x_1, \x_2, \cdots, \x_N\}$) . 
    \item $\H$ 是一个任务特定的希尔伯特空间, 其元素称为\textit{类补丁向量}. 
    \item 补丁 $\x$、特征 $\v$ 和核 $\w$ 都是类补丁向量. 
    \item 图像 $\z = (\x_1, \x_2, \cdots ,\x_P)^T \in \H^P$ 称为 $P$-补丁数据. 
    \item $P$ 是图像中补丁的数量. 
    \item $d$ 是 $\H$ 的维数, 也称为补丁/核的大小. 
    \item $C = c_2 d$ 是核的数量.  
    \item $\V := \{ \v_1, \v_2,\cdots, \v_d \}$ 是所有特征的集合. 
    \item $\V_{nor}$ 和 $\V_{aux}$ 分别是正常特征和辅助特征的集合. 
    \begin{itemize}
        \item 在理论分析中, $\V_{nor} = \{ \v_1, \v_2\}$ 且 $\V_{aux} = \{ \v_3, \v_4, \cdots , \v_d  \} $. 
        \item 我们在合成实验中探索了更多组合. 
    \end{itemize}
    \item 补丁 $\x_i$ \textit{包含}特征 $\v_k$ 意味着 $\x_i = \v_k + \rho$, 其中 $\rho$ 是小的随机噪声 (参见 \ref{sec-data-assumption}) . 
    $\epsilon$ 控制 $\rho$ 的大小. 
    \item $i \in [P]$, $j \in [C]$ 和 $k \in [d]$ 分别是补丁 ($\x_i$) 、核 ($\w_j$) 和特征 ($\v_k$) 的互斥指标. 
    \item $\phi = \phi_d \circ \phi_e$ 是在 \ref{sec-network-structure} 中定义的两层卷积自编码器. 
    \item $\sigma$ 是在 \ref{eq-def-smoothed-relu} 中定义的平滑 ReLU 函数, 由 $c_0$ 参数化. 
    \item $\ell(\z;\w)$ 和 $\risk_N(\w)$ 分别表示重构损失和误差. 
    \item $\sigma_0$ 控制 $\w$ 的 (零均值高斯) 随机初始化的大小. 
    \item $T$ 是最大迭代次数. 
    \item 给定 $t \in [T]$ 和 $k \in [d]$, $\smost^t(k)$ 表示特征 $\v_k$ 在第 $t$ 次迭代时的锥集. 
    \item $f_r = f_r(d,t)$, $f_h = f_h(d)$ 和常数 $c_1$ 控制锥的半径和高度. 
    \item $p_0 = p_0(d)$, $p_1 = p_1(d)$ 和 $p_2 = p_2(d)$ 控制概率界. 
    \begin{itemize}
        \item \ref{eq-lemma-size-of-wj-p0} 指定了 $p_0$ 的表达式. 
        引理 \ref{lemma-size-of-wj-in-paper} 中 $\|\w_j\|$ 的上界以至少 $1 - p_0$ 的概率成立. 
        \item $p_1$ (在 \ref{eq-def-p1} 中指定) 和 $p_2$ (在 \ref{eq-def-p2} 中指定) 控制 \ref{theorem-initial-size-in-paper} 中 $|\smost^{0}(k)|$ (锥集的初始大小) 的双侧界. 
    \end{itemize}
    
\end{itemize}

\section{参数分析}\label{sec-parameters}
选择一组合适的参数可以显著简化特征学习的分析, 特别是那些涉及紧不等式的证明. 
然而, 找到这样的合适参数极其费力, 因为每次尝试新的参数集都需要重新证明大部分结果. 
此外, 我们认为固定参数可能会使读者困惑, 并损害分析的整体直观性. 
在特征学习的分析中, 参数的选择不是唯一的； 
多种不同的选择可能导致相似的结果. 
因此, 固定参数的选择可能会掩盖不同变量之间的实际关系. 
基于上述原因, 本文指定参数的约束条件而不是固定它们. 

\subsection{推荐参数}
补丁维数 $d$、补丁数量 $P$ 以及假设 \ref{assumption-normal-data} 中的概率 $\beta_k$ ($k\in [d]$) 是任务特定的常数. 
为方便读者, 可取 $P = 10$, $d = 20$, 且对所有 $k > 2$ 有 $\beta_k < 0.5$. 

本文中的其他参数可分为三组: 
\begin{enumerate}
    \item 第一组是常数, 包括 $c_0$, $c_1$ 和 $c_2$. 建议取 $c_0 = 0.01$, $c_1 = 0.5$ 和 $c_2 = 1.2$. 
    \item 第二组包括依赖于维数 $d$ 的参数. 
    为简单起见, 我们令 $\|\v_k\|=1$. 
    为确保分析不失一般性, {噪声阈值} $\epsilon$ 和随机初始化的方差 $\sigma_0$ 应相应重新缩放. 
    为便于理解, 我们建议读者在本文中考虑 $\sigma_0 \leq d^{-2}$ 和 $\epsilon = \sigma_0$, 例如 $\epsilon = \sigma = 0.001$. 
    \item 第三组包括依赖于 $t$ 和 $d$ 的参数, 主要是与迭代相关的变量. 约束条件在本节其余部分讨论. 
\end{enumerate}

我们假设学习率 $\eta_t $ 以 $\eta_t = 0.2 t^{-1}$ 的速度随 $t$ 增长而衰减. 
为完整性, 我们令 $\eta_0 = 0$ (即, 核在第一次迭代中保持不变) . 
我们分析结果的证明对参数施加了几个约束. 
为确保引理 \ref{lemma-smost-delta-ij-in-paper} 成立, 我们假设
\begin{equation}
    (c_1 - 2\sqrt{2\epsilon}) f_h(d) = (0.5 - 2\sqrt{0.002}) f_h(d) > f_r(d,t),
\end{equation}
这是一个弱条件, 因为 $f_r(d,t) =  o(t) \cdot f_h(d)$ 且 $\epsilon \leq d^{-2}$. 
\ref{theorem-Stk-subset-st+1k} 的证明要求 $c_1$, $c_0$ 和 $\epsilon$ 满足
\begin{equation}
    c_1 - \sqrt{2\epsilon} > c_0,
\end{equation}
其中 $c_0$ 是 \ref{eq-def-smoothed-relu} 中平滑 ReLU 函数的阈值. 
记 $\tilde{\epsilon} = 3\sqrt{2\epsilon} + 4\epsilon 
\approx 0.138$, $\ \hat{\epsilon} = 1 + \sqrt{2\epsilon} \approx 1.045$, 以及 $\ \tilde{c}_1 = c_1 - \sqrt{2\epsilon} \approx 0.455$.  
作为我们分析的主要结果, \ref{theorem-Stk-subset-st+1k} 要求
\begin{equation}\label{eq-require-upper-bound}
N \beta_k \left( 
        % 2\tilde{c}_1\hat{\epsilon} - (1+\sqrt{2\epsilon})(\tilde{c}_1 + \hat{\epsilon})(\wbound)^2|\S_n(i_n)|
        2\tilde{c}_1\hat{\epsilon} - (1+\sqrt{2\epsilon})(\tilde{c}_1 + \hat{\epsilon})(\wbound)^2 (c_2 - 1) d
        \right) \geq \frac{5c_1}{2}.
\end{equation}
将推荐参数代入 \ref{eq-require-upper-bound}, 我们有
\begin{equation}
    N \beta_k \left( 
        % 2\tilde{c}_1\hat{\epsilon} - (1+\sqrt{2\epsilon})(\tilde{c}_1 + \hat{\epsilon})(\wbound)^2|\S_n(i_n)|
        2 \times 0.455 \times 1.045 - 1.045 \times 1.5 \times 10^{-6} \times 4 \times (t f_h)^2
        \right) \geq 1.25,
\end{equation}
即, 对 $N$, $\beta_k$, $t$ 和 $f_h$ 的第一个约束是
\begin{equation}
    N \beta_k \left( 0.951 - 6.27 \times 10^{-6} (tf_h)^2 \right) \geq 1.25.
\end{equation}
来自 \ref{theorem-Stk-subset-st+1k} 的另一个约束是
\begin{equation}\label{eq-require-lower-bound}
\begin{aligned}
    &0.4 N \beta_{k^{\prime}} \sigma_0 \left(
    f_h \cdot (6\epsilon + \sqrt{2\epsilon} + \sqrt{2\epsilon} \wbound (c_2 - 1)d ) + f_r \cdot \left(\wbound \cdot (c_2 - 1) d + \sqrt{2\epsilon}\right)
    \right) \\
    \leq &(t+1) \sigma_0 f_r(t+1, d).
\end{aligned}
\end{equation}
令 $T = N = 100$, $f_h(d) = 0.2\sqrt{d}$, 以及 $f_r(d, t) = f_h \cdot (t+1)^{-0.5}$. 
\ref{eq-require-lower-bound} 变为
\begin{equation}
    40 \beta_{k^{\prime}} \left( f_h \cdot (0.051 + 0.179 \times t \sigma_0 f_h) + f_r (4t \sigma_0 f_h + 0.045) \right) \leq (t+1) f_r(t+1, d)
\end{equation}
不难验证推荐参数的有效性. 